\section{Trust and Its Cost}
\label{sec:trust}

The design rests on one claim: attestation reduces a hostile owner to an unreliable one. This section states the claim, shows how a machine earns trust and a tenant's key, and defines the exposure that keys create.

\subsection{Threat model}
\label{sec:threat}
We trust the CPU, its TPM and its memory protection; the measured appliance image; each tenant's own key service; and the provider's control plane, for availability and accounting only (Table~\ref{tab:threat}).
Out of scope are laboratory-grade physical attacks such as voltage glitching of firmware TPMs~\cite{jacob2023faultpm}, flaws in the appliance, side channels and traffic analysis, and failure of the control plane. An owner who can read memory while the appliance runs (cold-boot or bus probing) is outside the claim below unless the host encrypts memory. We express that residual risk only as a policy number, the \emph{assurance} of Section~\ref{sec:exposure}.

\begin{table}[t]
\caption{Adversaries, what neutralises them, and what is left.}
\label{tab:threat}
\centering\small
\setlength{\tabcolsep}{4pt}
\begin{tabular}{@{}P{2.0cm}P{4.2cm}P{4.2cm}P{2.5cm}@{}}
\toprule
Adversary & Capability & Neutralised by & What remains \\
\midrule
Volunteer owner & controls hardware, power, network; may boot modified software, fake a TPM, replay, delay, drop & measured boot and attestation; manufacturer-certified keys; locked image with no login; disk sealed to the measurement & \textbf{omission and timing} \\
Network attacker & watches and changes traffic & encryption ending inside attested parties & traffic metadata \\
Other tenants & malicious function code & microVM isolation & side channels \\
\bottomrule
\end{tabular}
\end{table}

\textbf{Property 1 (threat reduction).} Under these assumptions, every behaviour of a worker or scheduler host that other components can observe, and that software can reach, is a crash, an omission or a delay.
The owner cannot run code other than the measured image, because it would fail attestation and receive no certificate and no key. The owner cannot read or change the image's memory or sealed state, and results and receipts are signed with a key that never leaves the image. What is left is power, the network path and the timing of when to stop answering.
So the mechanisms that churn already forces on us also cover malice, and redundancy is no longer needed for integrity.

Other mechanisms fall short. Memory encryption alone attests nothing and leaves plaintext visible to the operating system. Running every task several times and voting triples cost. Verifiable computation is far too slow for functions of milliseconds. Confidential VMs isolate best but do not exist on consumer hardware, so we treat them as one substrate among several (Section~\ref{sec:substrates}). A blockchain would protect against no adversary here, because a receipt signed by an attested key can already be checked by anyone (Section~\ref{sec:receipts}).

\subsection{Becoming an attested worker}
The appliance is a signed kernel image with a read-only, integrity-checked root file system, the worker agent and the scheduler code, with no owner login. The partition that holds caches is encrypted with a key sealed to the expected measurement. A session starts when the owner reboots into the image and ends on reboot, power loss, a drain request or failure.

Figure~\ref{fig:attest} shows the exchange, which follows Keylime-style attestation~\cite{schear2016keylime}. The agent creates a \emph{session key pair} inside the image and binds its public half into the quote. The verifier checks the quote and certificate chain, compares the boot record with reference values and, if all is well, issues a \emph{session certificate} with the machine's identity, the session public key, an assurance score and an expiry.
Every channel in the system uses this certificate, so no untrusted operating system sits on any path.
Two properties matter for scheduling. \emph{Trust is perishable:} the certificate expires when the session ends or one validity period after the last good quote, is extended by periodic re-quotes, and is revoked if one fails. And \emph{identity costs a machine:} a worker's identity is the hash of its TPM's endorsement key, so many identities cost real hardware~\cite{douceur2002sybil}.

\begin{figure}[t]
\centering
\resizebox{\columnwidth}{!}{\begin{tikzpicture}[x=1mm,y=1mm]
  \def\seqh{104}
  \seqactor{9}{ag}{Worker\\agent}
  \seqactor{33}{tp}{TPM}
  \seqactor{62}{vf}{Verifier}
  \seqactor{86}{kb}{Tenant\\KBS}
  \seqmsg{-11}{vf}{ag}{\circled{1} nonce}{rsp}
  \seqmsg{-17}{ag}{tp}{\circled{2} quote(PCRs, nonce)}{msg}
  \seqmsg{-22}{tp}{ag}{signed quote}{rsp}
  \seqmsg{-28}{ag}{vf}{\circled{3} ATTEST(quote, event log, session key)}{msg}
  \seqnote{62}{-34.5}{verify signature and EK chain; replay event\\log against reference values}
  \seqmsg{-42}{vf}{ag}{\circled{4} SESSION\_CERT(id, key, assurance, expiry)}{msg}
  \seqmsg{-50}{ag}{kb}{\circled{5} KEY\_REQUEST(cert [+ fresh quote])}{msg}
  \seqnote{86}{-56.5}{check cert and measurement\\version against tenant policy}
  \seqmsg{-63}{kb}{ag}{key wrapped to session key}{rsp}
  \seqnote{9}{-70}{keep key in memory only;\\decrypt tenant bundles}
  \draw[dashed, gray] (0,-76) rectangle (96,-96);
  \node[font=\tinyx, anchor=north west, inner sep=1pt] at (0,-76) {loop every $T_{\text{reattest}}$};
  \seqmsg{-82}{vf}{ag}{nonce}{rsp}
  \seqmsg{-87}{ag}{vf}{fresh quote $+$ log delta}{msg}
  \seqmsg{-93}{vf}{ag}{on failure: revoke cert, notify routers and scheduler}{msg}
\end{tikzpicture}
}
\caption{Attestation and key provisioning for one worker. The session key is bound into the quote, so the certificate belongs to this appliance instance only.}
\label{fig:attest}
\end{figure}

\subsection{Getting a tenant's key}
\label{sec:provisioning}
Attestation shows that a machine runs the appliance, not that it may see a particular tenant's data. The tenant's key service holds a policy, for instance ``accept appliance versions at or above this one, optionally with a fresh quote''.
Making worker $j$ \emph{provisioned} for tenant $t$ takes one round trip: the worker sends its certificate, the key service checks the policy and returns the key wrapped to the session public key, and the worker keeps it in memory only. This trip, plus decrypting bundles, is the provisioning term in the latency.
Code bundles are encrypted with the tenant key, so any worker may store and forward them but only provisioned workers can use them. Snapshots contain decrypted memory, so they are written only by provisioned workers, shared only with other provisioned workers of the same snapshot class, and counted in exposure.

Whom the tenant must trust for confidentiality depends on the policy. If the key service accepts any certificate the verifier issued for an approved measurement, the verifier is trusted for confidentiality. If the policy demands a fresh quote, the key service checks it itself and does not depend on the verifier, at the price of one more quote.
We assume the verifier acts as certificate authority, with its root distributed in the image and to key services. We do not evaluate this infrastructure, and the simulator charges one fixed calibrated cost for the policy check.

\subsection{Exposure and assurance}
\label{sec:exposure}
Every provisioned machine is a place from which a tenant's secret could leak if our trust assumptions fail. We define the \emph{key exposure} of tenant $t$ as
\begin{equation}
E_t = \sum_{j\in \mathrm{KH}(t)} (1-a_j),
\label{eq:exposure}
\end{equation}
where $\mathrm{KH}(t)$ is the set of machines holding $t$'s key or decrypted artifacts and $a_j\in[0,1]$ is a machine's \emph{assurance}: a base value for its trust substrate times a factor for firmware policy, fixed when the certificate is issued.
Exposure is a transparent proxy for attack surface, not a probability of compromise, and base assurances are policy parameters, not measured truths. We therefore report latency, exposure and robustness together as a trade-off front, never as one weighted number. Tenants set a sensitivity that weights exposure and an assurance floor below which a machine never gets their key.

\subsection{Substrates}
\label{sec:substrates}
A \emph{substrate profile} is data that tells one generic lifecycle engine how long trust takes, how long it lasts, what each start-up level costs and what assurance the host carries, so the scheduler is independent of how trust is established.
We model five. The main one, the TPM-attested appliance on consumer hardware, uses costs we measured (Section~\ref{sec:calib}). Two confidential-VM variants for server CPUs use ranges from the literature~\cite{amd2020snp,misono2025cvm,holmes2024severifast,li2023bifrost}. One models a measured hypervisor under the owner's operating system. The last, \texttt{none}, is the unattested reference that baselines use.
